\documentclass[10pt,a4paper]{article}
\usepackage{fontspec,amsmath,amssymb,graphicx}
\usepackage[margin=18mm]{geometry}
\setmainfont{Times New Roman}
\setlength{\parindent}{0pt}
\setlength{\parskip}{5pt}
\allowdisplaybreaks
\sloppy
\emergencystretch=2em
\begin{document}
\section*{VD\_4\_2\_2\_01 — Lời giải bổ sung}
\subsection*{Kiến thức và trình tự giải}
\textbf{Bước 1 — Viết toán tử và biên parabol.} Đặt $P=\partial_t-a(x,t)\partial_{xx}-b(x,t)\partial_x-c(x,t)$ với $a\ge0$, $c\le0$. Biên parabol gồm mặt $t=0$ và các biên không gian; mặt $t=T$ vẫn thuộc miền được ước lượng. Với phương trình nhiệt, $a=K>0$, $b=c=0$.

\textbf{Bước 2 — Loại cực đại dương bằng dấu đạo hàm.} Tại cực đại dương nội miền, $w_x=0$, $w_{xx}\le0$, $w_t=0$; nếu cực đại trên mặt thời gian trên thì đạo hàm một phía cho $w_t\ge0$. Suy $Pw\ge0$, mâu thuẫn $Pw<0$. Với $Pw\le0$, xét $w_\varepsilon=w-\varepsilon e^{\gamma t}$, $\gamma>0$: $Pw_\varepsilon=Pw-\varepsilon(\gamma-c)e^{\gamma t}<0$. Chứng minh $w_\varepsilon\le0$ rồi cho $\varepsilon\downarrow0$.

\textbf{Bước 3 — So sánh hai hàm bằng hiệu.} Để chứng minh $v\le u\le V$, kiểm tra $Pv\le Pu\le PV$ và thứ tự của dữ liệu đầu, biên. Áp dụng bước 2 cho $v-u$ và $u-V$. Nếu có nguồn dương, hằng số không nhất thiết là nghiệm trên, nên phải tính $PV-Pu$; không dùng cực đại của dữ liệu để chặn một bài có nguồn mà chưa kiểm tra.

\textbf{Bước 4 — Xử lý Neumann, Robin và tuần hoàn.} Neumann thuần nhất có thể phản xạ chẵn qua đầu mút; điểm biên trở thành điểm trong. Robin viết theo pháp tuyến ngoài $\partial_\nu u+\sigma u=0$, $\sigma\ge0$. Tại cực đại dương không hằng, bổ đề Hopf cho $\partial_\nu u>0$, trái với $\partial_\nu u=-\sigma u\le0$. Với Robin phải thêm mức $0$ vào cận hằng nếu dữ liệu mang một dấu. Biên tuần hoàn ghép hai đầu thành đường tròn; không còn biên không gian.

\textbf{Bước 5 — Miền vô hạn và nguồn bị chặn.} Trên toàn trục/nửa trục, cần nghiệm bị chặn trên mỗi dải thời gian (hoặc điều kiện tăng trưởng đủ dùng). Dùng rào $\varepsilon((x-x_0)^2+2Kt)$ trên đoạn lớn, áp dụng cực đại trên miền hữu hạn rồi cho $\varepsilon\downarrow0$; cách này loại ảnh hưởng từ vô cực. Nếu $|Pu|\le N$ và $|u|\le M$ trên biên parabol, đặt $V=M+Nt$. Vì $PV=N-cV\ge N$, suy $|u|\le M+Nt\le M+NT$. $M$ ở đây là cận của $|u|$, không phải cực đại có dấu.
\subsection*{Lời giải áp dụng cho bài}
\textbf{Bước 1. Đọc đúng dữ liệu.} Đề là $u_t=u_{xx}$ trên $0<x<1$, biên $u(0,t)=u(1,t)=0$ và $f=\sin^2(\pi x)$. Ký hiệu trong bản scan là $\pi$, không phải tần số nguyên $n$; không cần đặt lại miền hay giới hạn $n$.
\textbf{Bước 2. Chọn hai nghiệm so sánh.} Đặt $v=0$, $V=e^{-\pi^2t}\sin\pi x$. Tính $V_t=-\pi^2V=V_{xx}$, $V(0,t)=V(1,t)=0$. Trên $[0,1]$, $0\le\sin\pi x\le1$, nên $0\le\sin^2\pi x\le\sin\pi x$.
\textbf{Bước 3. Áp dụng cực đại cho hai hiệu.} $v-u=-u$ giải phương trình nhiệt, dữ liệu đầu và biên không dương, nên $-u\le0$. Tương tự $u-V$ có dữ liệu đầu $\sin\pi x(\sin\pi x-1)\le0$ và biên $0$, nên $u-V\le0$. Kết luận
\[0\le u(x,t)\le e^{-\pi^2t}\sin\pi x,\quad 0\le x\le1,\ t\ge0.\]
Nếu cần chứng minh cực đại yếu, xét mỗi hiệu trừ $\varepsilon e^t$; toán tử $\partial_t-\partial_{xx}$ cho giá trị âm nghiêm ngặt, loại cực đại dương, rồi cho $\varepsilon\downarrow0$.\end{document}